\section{Every directional Laurent coefficient of fully symmetrized depth}
\label{sec:symmetric}

At depth greater than two, resolving one ordered multiple-zeta germ
requires additional information. Complete symmetrization nevertheless
admits an exact answer at every depth and every Laurent order. The
algebraic input is Hoffman's classical symmetric-sum principle
\cite{hoffman}; the present use records independent perturbation
directions and the resulting Stieltjes coefficients explicitly.

For $a>0$ and $d\ge1$ set
\[
 Z_d(s_1,\ldots,s_d;a)
   =\sum_{n_1>\cdots>n_d\ge0}\prod_{i=1}^d(n_i+a)^{-s_i},
\]
initially in its domain of absolute convergence. We sum over all
\emph{labelled} permutations,
\[
 \Sym Z_d(\boldsymbol s;a)
   =\sum_{\sigma\in S_d}Z_d(s_{\sigma(1)},\ldots,s_{\sigma(d)};a).
\]
Repeated exponent values still occur with their permutation
multiplicity. Let $\Pi_d$ denote the set partitions of $\{1,\ldots,d\}$.

\begin{proposition}[Classical partition identity]\label{prop:partition}
As an identity of meromorphic functions,
\begin{equation}\label{eq:partition}
 \Sym Z_d(\boldsymbol s;a)
 =\sum_{\pi\in\Pi_d}(-1)^{d-|\pi|}
      \prod_{B\in\pi}\left[(|B|-1)!\,
                \zeta\left(\sum_{i\in B}s_i,a\right)\right].
\end{equation}
\end{proposition}
\begin{proof}
First truncate all indices to $0\le n<N$. The symmetrized ordered
sum is the sum over all labelled tuples with distinct indices.
Inclusion--exclusion on equality patterns is M\"obius inversion in the
partition lattice. Its coefficient for a block $B$ is
$(-1)^{|B|-1}(|B|-1)!$. Each unrestricted equality block contributes
$\sum_{n<N}(n+a)^{-\sum_{i\in B}s_i}$. This proves the finite identity.
Letting $N\to\infty$ in the absolute-convergence region gives
\eqref{eq:partition}; meromorphic continuation gives the asserted
identity. Alternatively, the right side itself supplies the
continuation of the symmetrized function used below.
\end{proof}

\begin{theorem}[All-depth, all-order directional coefficient algorithm]
\label{thm:symcoeff}
Let $c_1,\ldots,c_d\in\C\setminus\{0\}$ and put
$s_i=1+c_i\eps$. For a block $B$, write $b=|B|$ and
$c_B=\sum_{i\in B}c_i$. Define the coefficient array
\begin{align}
 A_{\{i\}}(-1)&=\frac1{c_i},&
 A_{\{i\}}(k)&=\frac{(-1)^k\gamma_k(a)c_i^k}{k!}\quad(k\ge0),
 \label{eq:singletoncoeff}\\
 A_B(k)&=\frac{\zeta^{(k)}(b,a)c_B^k}{k!}\quad(k\ge0, b\ge2).
 \label{eq:blockcoeff}
\end{align}
All other entries are zero. Then, for every integer $q\ge-d$,
\begin{equation}\label{eq:symcoeff}
 \boxed{
 [\eps^q]\Sym Z_d(1+c_1\eps,\ldots,1+c_d\eps;a)
 =\sum_{\pi\in\Pi_d}(-1)^{d-|\pi|}
       \prod_{B\in\pi}(|B|-1)!
       \sum_{\substack{(k_B)_{B\in\pi}\\\sum_Bk_B=q}}
                     \prod_{B\in\pi}A_B(k_B).}
\end{equation}
For each fixed $d,q$ the displayed expression is finite. Vanishing
block sums $c_B$ with $|B|\ge2$ cause no singularity. The constant
coefficient uses only $\gamma_0(a),\ldots,\gamma_{d-1}(a)$ and
finitely many derivatives $\zeta^{(j)}(b,a)$ at integers $2\le b\le d$.
\end{theorem}
\begin{proof}
For a singleton, the zeta factor in \eqref{eq:partition} has exactly
the Laurent coefficients \eqref{eq:singletoncoeff}. For a block of
size at least two it is regular at $\eps=0$, with Taylor coefficients
\eqref{eq:blockcoeff}. Multiplication proves \eqref{eq:symcoeff}.
There are finitely many partitions, each index has lower bound $-1$
or $0$, and their sum is prescribed; hence only finitely many tuples
can contribute. For $q=0$, a singleton index is at most $d-1$,
which proves the Stieltjes-index assertion.
\end{proof}

This is an explicit finite algorithm: enumerate set partitions,
generate each block's coefficient list to the required degree, and
multiply the short Laurent polynomials. No nested infinite sums or
unspecified finite parts remain in the output.

\subsection{Depth two and depth three in closed form}

For depth two, the constant coefficient is
\begin{equation}\label{eq:sym2}
 \FP\Sym Z_2
 =\gamma_0(a)^2-\zeta(2,a)
        -\gamma_1(a)\left(\frac{c_1}{c_2}+\frac{c_2}{c_1}\right),
\end{equation}
in agreement with the sum of the two ordered values in
\eqref{eq:direction-FP} whenever both ordered restrictions are
transverse. Formula \eqref{eq:sym2} itself also applies when
$c_1+c_2=0$: the polar divisor of each ordered summand has canceled
in the symmetrized meromorphic germ before restriction.

For depth three let $c_\Sigma=c_1+c_2+c_3$. Direct expansion gives
\begin{align}\label{eq:sym3}
 \FP\Sym Z_3={}&\gamma_0(a)^3
   -\gamma_0(a)\gamma_1(a)\sum_{i\ne j}\frac{c_j}{c_i}
   +\frac{\gamma_2(a)}2\sum_{i=1}^3
                         \frac{c_i^2}{\prod_{j\ne i}c_j}\notag\\
 &-3\gamma_0(a)\zeta(2,a)
   -\zeta'(2,a)\sum_{i=1}^3\frac{c_\Sigma-c_i}{c_i}
   +2\zeta(3,a).
\end{align}
At $c_1=c_2=c_3=1$ this becomes
\begin{equation}\label{eq:sym3diagonal}
 \gamma_0(a)^3-6\gamma_0(a)\gamma_1(a)+\frac32\gamma_2(a)
 -3\gamma_0(a)\zeta(2,a)-6\zeta'(2,a)+2\zeta(3,a).
\end{equation}
The corresponding single diagonal depth-three finite part is one
sixth of \eqref{eq:sym3diagonal}; the factor $3!$ is required by our
labelled convention.

\subsection{Why coefficient extraction is not automatically a renormalization}

The Laurent constant of a product is generally different from the
product of the Laurent constants. For example,
\begin{equation}\label{eq:FPnonmult}
 \FP_{\eps=0}\bigl\{\zeta(1+c\eps,a)\zeta(1+d\eps,a)\bigr\}
 =\gamma_0(a)^2-\gamma_1(a)(c/d+d/c),
\end{equation}
whereas the product of the two individual constant coefficients is
$\gamma_0(a)^2$. Therefore one cannot assert that a collection of
directional finite parts preserves the stuffle algebra simply because
the meromorphic functions satisfy stuffle identities. Algebra-preserving
renormalizations require a compatible subtraction prescription; this is
the subject of, for example, Guo and Zhang \cite{guozhang}.

Theorems~\ref{thm:doublegerm} and \ref{thm:symcoeff} settle two precise
parts of the incoming questions: every ordered depth-two transverse
direction, and every fully symmetrized direction at arbitrary depth.
They do not give each individual ordered germ at depth three and higher,
nor do they identify a canonical renormalization independent of a chosen
path and subtraction convention.

